\begin{trapbox}
	\begin{description}[style=nextline,leftmargin=2.8em,labelsep=0.5em,itemsep=0.35em]
		\item[\textbf{\textcolor{CTrap}{易错点一（忽略定义域）}}] 未检查真数 $b>0$ 就应用符号法则。
		\item[\textbf{\textcolor{CTrap}{正确理解（定义域优先）：}}] 使用 $\log_a b$ 的符号法则前，必须先确保 $b>0$，且 $a>0,\ a\neq1$。
		\item[\textbf{\textcolor{CTrap}{错因分析（跳过验真）：}}] 对数运算要求真数严格正，跳过此步骤直接计算 $(a-1)(b-1)$ 可能得到错误符号（如真数为负时无意义）。
	\end{description}

	\medskip
	\begin{description}[style=nextline,leftmargin=2.8em,labelsep=0.5em,itemsep=0.35em]
		\item[\textbf{\textcolor{CTrap}{易错点二（片面记忆）}}] 错误认为底数 $a>1$ 时 $\log_a b$ 一定为正，或底数 $0<a<1$ 时一定为负。
		\item[\textbf{\textcolor{CTrap}{正确理解（乘积判号）：}}] $\log_a b$ 的符号取决于 $(a-1)(b-1)$ 的符号，即底数与真数相对于 $1$ 的大小关系共同决定。
		\item[\textbf{\textcolor{CTrap}{错因分析（忽略真数）：}}] 只考虑底数，忽略了真数 $b$ 的位置。例如 $\log_2 0.5 = -1 < 0$，底数虽大于 $1$，但真数小于 $1$，结果反而为负。
	\end{description}

	\medskip
	\begin{description}[style=nextline,leftmargin=2.8em,labelsep=0.5em,itemsep=0.35em]
		\item[\textbf{\textcolor{CTrap}{易错点三（忽视零界）}}] 将 $b=1$ 或 $a=1$ 的情况代入符号法则，认为 $\log_a 1 >0$ 或 $<0$。
		\item[\textbf{\textcolor{CTrap}{正确理解（等号排除）：}}] 当 $b=1$ 时，$\log_a b = 0$，不满足 $>0$ 或 $<0$；当 $a=1$ 时对数无定义，也应排除。
		\item[\textbf{\textcolor{CTrap}{错因分析（机械套用）：}}] 直接套用 $(a-1)(b-1)>0$ 的条件，得到 $b=1$ 时 $(a-1)(0)=0$，不属于正或负，导致误判符号。
	\end{description}
\end{trapbox}
